\chapter{Growth \& Moats}\label{ch:growth}

% Scope: the Ansoff matrix, the moat taxonomy (switching costs, network effects, scale,
%   intangibles), network-effect mechanics, growth accounting.
% Cross-link: network effects are the structural moat; the operational flywheel that generates
%   them is in \cref{ch:digital-funnel-and-crm}. Product-market expansion <-> STP
%   (\cref{ch:segmentation-targeting-positioning}).

A positive advantage spread (\cref{eq:advantage-spread}) answers whether you have a moat;
it does not say how you grew into it or how you protect and widen it. Growth strategy is
the set of decisions about which markets and products to enter next; moat strategy is the
set of structural choices that make the resulting position defensible. Both operate through
the same financial bottleneck: every growth move must generate cohort NPV that exceeds its
cost.

\section{The Ansoff Matrix}\label{sec:ansoff-matrix}
% complexity: 3

Which of the four growth vectors --- sell more to existing customers, launch new products,
enter new markets, or diversify --- offers the best return per unit of execution risk? Growth
always involves a bet on something unfamiliar. The Ansoff matrix makes the bet explicit: you
are betting on a known product in a known market (low risk), a new product in a known market,
a known product in a new market, or a new product in a new market (high risk). The risk rises
diagonally, and so does the required return.

\textcite{kotler2021} present the matrix as a \(2\times2\) decision structure across
product familiarity (existing / new) and market familiarity (existing / new):
\begin{description}
  \item[Market penetration] Existing product, existing market. Grow share via price,
        promotion, or distribution. Lowest risk; highest resource efficiency; limited
        ceiling in a mature market.
  \item[Product development] New product, existing market. Extend the platform for
        customers the firm already understands. Moderate risk (product uncertainty, but
        customer knowledge retained).
  \item[Market development] Existing product, new market. Replicate the value proposition
        in a new geography, vertical, or segment. Moderate risk (product known, customer
        unknown). Requires a new STP decision (\cref{ch:segmentation-targeting-positioning}).
  \item[Diversification] New product, new market. The highest-risk quadrant; justified only
        when core markets are saturated or declining, and usually executed via M\&A rather
        than organic build.
\end{description}

``Market'' and ``product'' are definitional choices, not objective facts. The same move can
be labelled penetration (``we're just adding features'') or product development (``this is
a new product line'') depending on framing. The useful version of the matrix forces a
concrete judgment: does the new offering require capabilities or go-to-market motions the
firm does not currently have? If yes, it is a new-product or new-market move; if no, it is
penetration. Without that rigor the matrix becomes a vocabulary, not a tool.

\begin{example}
\emph{Decision: which growth vector should the SaaS firm prioritise in the next 18 months?}

Current state: \money{4000000} ARR, \num{200} new customers/month, \money{340000} MRR.
The four options evaluated against cohort NPV:

\begin{itemize}
  \item \textbf{Market penetration}: the funnel converts \qty{2.5}{\percent}
        (200/8000 qualified sessions). A/B testing landing pages and improving the
        trial-to-paid step (\cref{ch:digital-funnel-and-crm}) could lift conversion to
        \qty{3}{\percent} --- 240 paid/month --- at near-zero incremental cost.
        Highest NPV per dollar invested at current scale.
  \item \textbf{Product development}: a reporting add-on at \money{20}/month would raise
        ARPU from \money{50} to \money{70}, lifting LTV from \money{3150} to
        \(\money{70}\times 0.42/(0.11-0.03) = \money{4410}\). The constraint is
        whether existing customers want it (\cref{ch:segmentation-targeting-positioning}).
  \item \textbf{Market development}: entering the European SMB segment requires GDPR
        compliance, local payment methods, and a localised onboarding flow ---
        meaningful execution cost with uncertain CAC.
  \item \textbf{Diversification}: premature at \money{4000000} ARR; reserves of
        execution attention are the binding constraint.
\end{itemize}

Penetration and product development are the rational first moves; market development
is a 24-month option once unit economics are proven.
\end{example}

\begin{admonition}{Warning}
Using the Ansoff matrix as a menu rather than a sequencing tool. All four options are
always theoretically available; the matrix's value is in forcing an explicit ranking by
risk-adjusted return and execution capacity. Pursuing two high-risk quadrants simultaneously
dilutes management attention below the minimum effective dose in either.
\end{admonition}

The Ansoff matrix is the corporate-strategy analogue of the segmentation decision in
\cref{ch:segmentation-targeting-positioning}: both ask the same question (who are we serving
and with what?), one at the portfolio level and one at the product level. Corporate-strategy
research (\parencite{dranove7theconomics}) finds that \emph{related} diversification ---
moving along a single dimension of the matrix at a time --- generates superior returns to
unrelated diversification, consistent with the VRIN argument that advantage is resource-
specific and does not transfer across unrelated businesses.

The Ansoff sequencing should follow the moat taxonomy below: penetration first, product
development once the moat is established, market development once the product is proven.
\textcite{kotler2021} contextualise the matrix within modern growth strategy.

\section{The Moat Taxonomy}\label{sec:moat-taxonomy}
% complexity: 3

Which structural barrier will protect the firm's returns once a competitor has the resources
to match product quality and price? The moat must be built before it is needed. A moat is
any structural advantage that makes it irrational for a well-resourced competitor to attack
you directly. Unlike a product advantage, a moat is self-reinforcing: the more customers you
serve, the harder it is for a rival to dislodge them. Moats turn a temporary lead into a
sustained positive value spread (\cref{eq:advantage-spread}).

Four moat types sustain ROIC $>$ WACC (\cref{eq:roic}) over a strategic horizon
\parencite{porter1985,dranove7theconomics}:
\begin{description}
  \item[Switching costs] The direct and indirect cost a customer incurs when leaving for a
        rival. Includes data migration, retraining, workflow disruption, and integration
        rebuilds. Switching costs raise the effective price of leaving without raising the
        nominal price of staying --- they expand the price band within which the incumbent
        is safe from defection.
  \item[Network effects] Each additional user makes the product more valuable to every
        other user. This is the most powerful moat in technology markets (formalised in
        \cref{sec:network-effects}): it is demand-side scale, not cost-side scale, and
        it compounds automatically.
  \item[Cost-side scale] Declining unit costs at large scale --- spreading fixed cost,
        purchasing leverage, or learning-curve effects. Protects via price: incumbents can
        profitably undercut an entrant who has not yet reached minimum efficient scale.
  \item[Intangible assets] Patents, regulatory licences, or brand reputation that
        competitors cannot replicate quickly. Brand is the weakest intangible moat in
        B2B SaaS (buyers are professional and price-sensitive); patents require dense
        IP investment.
\end{description}

Moats erode. Switching costs collapse when a standard data-portability layer emerges (e.g.\
GDPR's right to data portability in SaaS). Network effects reverse when a platform reaches
diseconomies of scale (spam, congestion). Cost-side scale is undermined by a disruptive
technology that resets the cost structure. Treat a moat as a wasting asset that requires
reinvestment, not a permanent barrier.

\begin{example}
\emph{Decision: what is the SaaS firm's primary moat and how much is it worth?}

Primary moat: \emph{switching costs} from integration depth. A customer with 10 active
integrations faces an estimated \money{3000}--\money{6000} of migration cost
(engineering time, retraining, downtime risk) --- well above the annual subscription value
of \money{600}. This means the customer's effective switching cost exceeds one full year of
revenue, materially suppressing churn below the structural average.

Financial expression: if the switching-cost moat reduces annual churn from
\qty{10}{\percent} to \qty{7}{\percent}, the monthly net churn rate shifts from
\qty{0.87}{\percent} to \qty{0.60}{\percent} and LTV rises:
\[
  \Delta\mathrm{LTV} \approx \frac{\money{252}}{0.08} - \frac{\money{252}}{0.11}
                           = \money{3150} - \money{2291} \approx \money{860}.
\]
A \money{860} rise in LTV means the firm can spend \money{860} more per customer acquired
without violating the LTV:CAC floor (\cref{eq:ltv-cac}), which at 200 customers/month
unlocks \money{172000}/month of incremental CAC budget --- \qty{143}{\percent} of the
current ad spend.
\end{example}

\begin{admonition}{Warning}
Calling \emph{retention} a moat. Retention is the symptom; the moat is the structural
cause. A product that retains customers because it is excellent loses them the day a
competitor builds a better product. A product that retains customers because integrations
make switching painful retains them even when a better product exists. Invest in the cause,
not the symptom.
\end{admonition}

Moat analysis connects to the price-theory literature via the Lerner index
(\cref{eq:lerner}): a firm with switching costs can price above marginal cost by an amount
proportional to the switching cost divided by the contract horizon. The moat is, in
price-theory terms, a source of inelasticity: it reduces the own-price elasticity of demand
(\cref{eq:elasticity}) for the incumbent's product without changing the elasticity of the
overall category.

Switching costs are the SaaS firm's primary moat; network effects (\cref{sec:network-effects})
are the potential secondary moat if an integration marketplace is built. Together they
generate the sustained value spread that \cref{ch:advantage-and-disruption} defines.
\textcite{porter1985} and \textcite{dranove7theconomics} are the primary sources.

\section{Network-Effect Mechanics}\label{sec:network-effects}
% complexity: 4 -> reuse figure growth-loop

Does the product have network effects, how strong are they, and what does that imply for
the CAC budget and the growth trajectory? Network effects create a demand-side moat that
compounds: each new user raises the value of the product for all existing users, which makes
the product easier to sell to the next user, which lowers CAC over time. The structural
implication is that unit economics \emph{improve} with scale --- the opposite of most cost
structures.

Metcalfe's Law approximates the value of a network of \(n\) nodes as proportional to the
number of distinct pairs:
\begin{equation}\label{eq:metcalfe}
  V(n) \;\propto\; n(n-1) \;\approx\; n^2 \quad\text{for large }n.
\end{equation}
The quadratic scaling is the economic core: doubling users more than doubles value,
so value-per-user rises with network size. This means willingness to pay rises as the
network grows, without any change in the product itself --- the demand curve shifts outward
purely from the size of the installed base.

\begin{figure}[htbp]
  \centering
  \includegraphics[width=0.86\linewidth]{growth-loop}
  \caption{The self-reinforcing growth loop. New customers produce referrals, content, or
    data that attract further new customers. When the viral coefficient \(k > 1\)
    (\cref{eq:viral-k}), the loop is self-sustaining and CAC from paid channels falls as
    organic acquisition rises.}
  \label{fig:network-effects}
\end{figure}

\cref{eq:metcalfe} assumes all links are equally valuable, which is empirically false.
Real networks are power-law distributed: a small number of high-value nodes (power users,
enterprise accounts) contribute disproportionately to network value. The quadratic formula
overstates value in sparse, heterogeneous networks. It also conflates \emph{direct} network
effects (user benefits from other users directly) with \emph{indirect} network effects (user
benefits from complements attracted by a large user base). Direct effects are stronger but
rarer; indirect effects are more common in platform markets (\cref{ch:business-models}).

\begin{example}
\emph{Decision: if the SaaS firm builds an integration marketplace connecting vendors,
what is the potential value at \(n = 100\) integrations?}

From \cref{eq:metcalfe}, the number of distinct vendor pairs at \(n = 100\):
\[
  \frac{n(n-1)}{2} = \frac{100 \times 99}{2} = 4{,}950 \text{ pairs}.
\]
Each pair represents a potential data-sharing or workflow-automation opportunity that did
not exist at \(n = 50\) (which offered only \num{1225} pairs). The \emph{marginal}
value of the 100th integration is \(99\) new pairs --- 99 times the value of the first.
This is why integration marketplaces justify a disproportionate share of engineering
investment at early stages: the option value of the 100th node is embedded in the
first node's decision to join.

The growth loop (\cref{fig:network-effects}) captures the organisational expression: each
integration attracts more vendors, whose participation attracts more customers, whose
data improves the product, which attracts more integrations. The viral coefficient
(\cref{eq:viral-k}) is the demand-side instantiation of the same loop applied to users
rather than integrations. Both are channels to \cref{ch:digital-funnel-and-crm}.
\end{example}

\begin{admonition}{Warning}
Treating weak same-side effects as Metcalfe-scale network effects. A SaaS product used
independently by each customer (no shared data, no interactions between users) has zero
direct network effect regardless of user count. The relevant test is: does user \(j\)'s
experience change when user \(k\) joins? If not, it is not a network effect; it is a
product feature that happens to scale. Overstating network effects misleads capital
allocation.
\end{admonition}

Direct vs indirect network effects carry different implications for critical mass.
Direct effects require a minimum viable network of active users before value exceeds the
standalone product; indirect effects require a minimum viable set of complementors.
The chicken-and-egg problem in platform launch (\cref{ch:business-models}) is the
demand-side version of the critical-mass requirement in \cref{eq:metcalfe}.
\parencite{dranove7theconomics} develop the welfare and pricing implications.

Network effects are the strongest moat in technology markets but the hardest to build from
scratch. The viral coefficient (\cref{eq:viral-k}) in \cref{ch:digital-funnel-and-crm}
is the demand-side mechanism; the metrics for tracking it are in \cref{ch:metrics-catalog}.
\textcite{dranove7theconomics} provide the theoretical foundation.

\section{Growth Accounting}\label{sec:growth-accounting}
% complexity: 4 -> figure mrr-waterfall

Is the firm actually growing, or is new revenue simply replacing churned revenue? The
decomposition of net MRR growth tells you which channels to invest in and where the
leaks are. Aggregate revenue growth conflates four very different business activities. New
customer acquisition, expansion of existing accounts, contraction of downgraded accounts,
and churn of lost accounts all net to a single headline number that masks the underlying
dynamics. A firm with strong new sales and high churn looks identical in the headline to
a firm with modest new sales and strong retention, but they require completely different
strategic responses.

Let MRR at the start of a period be \(M_0\). The ending MRR is:
\begin{equation}\label{eq:mrr-growth}
  M_1 \;=\; M_0 \;+\; \Delta_{\text{new}}
               \;+\; \Delta_{\text{expansion}}
               \;-\; \Delta_{\text{contraction}}
               \;-\; \Delta_{\text{churn}},
\end{equation}
where each \(\Delta\) is a positive number. Net new MRR is
\(\Delta_{\text{new}} + \Delta_{\text{expansion}} - \Delta_{\text{contraction}} - \Delta_{\text{churn}}\).
Net Revenue Retention (NRR), derived in \cref{ch:business-models}, is:
\[
  \mathrm{NRR} = \frac{M_0 + \Delta_{\text{expansion}} - \Delta_{\text{contraction}} - \Delta_{\text{churn}}}{M_0},
\]
i.e.\ what revenue retention looks like \emph{excluding} new customers. NRR $>$ 100\%
means the existing base is expanding on net --- the firm could theoretically cease new
sales and still grow.

\begin{figure}[htbp]
  \centering
  \includegraphics[width=0.86\linewidth]{mrr-waterfall}
  \caption{MRR waterfall for the running SaaS example. Starting MRR of \money{340000}
    grows by new (\money{30000}) and expansion (\money{20000}), and shrinks by contraction
    (\money{5000}) and churn (\money{22000}), ending at \money{363000} --- a \qty{6.8}{\percent} net gain.}
  \label{fig:mrr-waterfall}
\end{figure}

\cref{eq:mrr-growth} is an identity --- it always holds by construction. The diagnostic
power comes from assessing whether each component is healthy \emph{relative to the
business model}. A high-velocity transactional product (e-commerce SaaS) is expected to
have higher gross churn and lower expansion; a seat-expansion enterprise product is expected
to have high expansion and low gross churn. Comparing absolute levels of churn without
adjusting for business-model type produces false benchmarks. Relatedly, the equation
treats all MRR as equally valuable; a churn dollar from a 36-month-tenure customer is more
concerning than one from a one-month-old account (\cref{eq:clv}).

\begin{example}
\emph{Decision: where should the team invest --- new acquisition or churn reduction?}

From the fact sheet: MRR = \money{340000}. Observed monthly movements:
Observed monthly movements (stated assumptions; plan mix and expansion push
\(\Delta_{\text{new}}\) above the simple \num{200}\(\times\)\money{50} product):
\[
  \begin{aligned}
    \Delta_{\text{new}} &= +\money{30000} \quad (\text{new customers across plan mix})\\
    \Delta_{\text{expansion}} &= +\money{20000} \quad (\text{seat adds, plan upgrades})\\
    \Delta_{\text{contraction}} &= -\money{5000} \quad (\text{downgrades})\\
    \Delta_{\text{churn}} &= -\money{22000} \quad (\approx\qty{6.5}{\percent}\text{ of }\money{340000})
  \end{aligned}
\]
Net new MRR: $\money{30000}+\money{20000}-\money{5000}-\money{22000}=+\money{23000}$
(+\qty{6.8}{\percent}/month).
\[
  \mathrm{NRR} = \frac{\money{340000}+\money{20000}-\money{5000}-\money{22000}}{\money{340000}}
               = \frac{\money{333000}}{\money{340000}} \approx \qty{97.9}{\percent}.
\]
NRR of \qty{97.9}{\percent} indicates the existing base is slightly contracting without
new sales --- the firm is net-retention negative. Even modest improvements to retention
compound dramatically: reducing churn from \money{22000} to \money{15000}/month
raises NRR to \qty{99.7}{\percent} and raises ending MRR by \money{7000}/month
without acquiring a single new customer. That \money{7000} compounds. The investment
priority is churn reduction --- specifically the VRIN resources (integrations) that
raise switching costs (\cref{sec:moat-taxonomy}).
\end{example}

\begin{admonition}{Warning}
Optimising \emph{gross} new MRR while ignoring churn. A firm adding \money{40000}/month
in new MRR while churning \money{40000}/month is running in place: NRR $=$ \qty{100}{\percent}
exactly, but the customer base has completely turned over in some finite horizon and the
CAC has been spent with zero compound return. \emph{Net} MRR growth and NRR are the
correct metrics; gross new MRR is a vanity number without the retention context.
\end{admonition}

The MRR waterfall is the operational version of the cohort-level LTV model in
\cref{ch:customer-economics}: \cref{eq:clv} discounts a single cohort's expected cash
flows; \cref{eq:mrr-growth} is the aggregate monthly manifestation of all overlapping
cohorts summed in calendar time. The metrics catalog (\cref{ch:metrics-catalog}) provides
benchmarks for each component. \textcite{ries2011} introduced the ``growth accounting''
framing for product metrics; \textcite{bendle2020marketing} extend it to the marketing
attribution context.

Growth accounting closes the loop between the growth-strategy choices (Ansoff) and the
moat investments (switching costs, network effects) by making the financial result of each
visible in real time. \cref{eq:mrr-growth} and NRR are the primary dashboard for
\cref{ch:metrics-catalog}. \textcite{ries2011} and \textcite{bendle2020marketing} are the
primary sources.
